\documentclass[10pt,a4paper]{amsart}
\usepackage{amsmath,amssymb,amsthm,graphicx}
% For commutative diagrams we use Xy-pic
\usepackage[all]{xy}

% We use 2cell for 2-commutative diagrams.
\xyoption{2cell}
% Theorem environments.
%
\theoremstyle{plain}
\newtheorem{theorem}[subsection]{Theorem}
\newtheorem{proposition}[subsection]{Proposition}
\newtheorem{lemma}[subsection]{Lemma}

\theoremstyle{definition}
\newtheorem{definition}[subsection]{Definition}
\newtheorem{example}[subsection]{Example}
\newtheorem{exercise}[subsection]{Exercise}
\newtheorem{situation}[subsection]{Situation}

\theoremstyle{remark}
\newtheorem{remark}[subsection]{Remark}
\newtheorem{remarks}[subsection]{Remarks}

\numberwithin{equation}{subsection}

% Macros
%
\def\lim{\mathop{\mathrm{lim}}\nolimits}
\def\colim{\mathop{\mathrm{colim}}\nolimits}
\def\Spec{\mathop{\mathrm{Spec}}}
\def\Hom{\mathop{\mathrm{Hom}}\nolimits}
\def\Ext{\mathop{\mathrm{Ext}}\nolimits}
\def\SheafHom{\mathop{\mathcal{H}\!\mathit{om}}\nolimits}
\def\SheafExt{\mathop{\mathcal{E}\!\mathit{xt}}\nolimits}
\def\Sch{\mathit{Sch}}
\def\Mor{\mathop{\mathrm{Mor}}\nolimits}
\def\Ob{\mathop{\mathrm{Ob}}\nolimits}
\def\Sh{\mathop{\mathit{Sh}}\nolimits}
\def\NL{\mathop{N\!L}\nolimits}
\def\CH{\mathop{\mathrm{CH}}\nolimits}
\def\proetale{{pro\text{-}\acute{e}tale}}
\def\etale{{\acute{e}tale}}
\def\QCoh{\mathit{QCoh}}
\def\Ker{\mathop{\mathrm{Ker}}}
\def\Im{\mathop{\mathrm{Im}}}
\def\Coker{\mathop{\mathrm{Coker}}}
\def\Coim{\mathop{\mathrm{Coim}}}

% Boxtimes
%
\DeclareMathSymbol{\boxtimes}{\mathbin}{AMSa}{"02}

%
% Macros for moduli stacks/spaces
%
\def\QCohstack{\mathcal{QC}\!\mathit{oh}}
\def\Cohstack{\mathcal{C}\!\mathit{oh}}
\def\Spacesstack{\mathcal{S}\!\mathit{paces}}
\def\Quotfunctor{\mathrm{Quot}}
\def\Hilbfunctor{\mathrm{Hilb}}
\def\Curvesstack{\mathcal{C}\!\mathit{urves}}
\def\Polarizedstack{\mathcal{P}\!\mathit{olarized}}
\def\Complexesstack{\mathcal{C}\!\mathit{omplexes}}
% \Pic is the operator that assigns to X its picard group, usage \Pic(X)
% \Picardstack_{X/B} denotes the Picard stack of X over B
% \Picardfunctor_{X/B} denotes the Picard functor of X over B
\def\Pic{\mathop{\mathrm{Pic}}\nolimits}
\def\Picardstack{\mathcal{P}\!\mathit{ic}}
\def\Picardfunctor{\mathrm{Pic}}
\def\Deformationcategory{\mathcal{D}\!\mathit{ef}}
\usepackage{cleveref}
\crefname{originalreference}{original reference}{original references}
\Crefname{originalreference}{Original reference}{Original references}
\makeatletter
\newlabel{lemma-restrict-to-open-coefficients}{{0GJE}{source}{}{sourcepointer-1}{}}
\newlabel{lemma-restrict-to-open-coefficients@cref}{{[originalreference][0][]0GJE}{[1][0][]source}{}{}{}}
\makeatother
\crefname{originalreference}{original reference}{original references}
\Crefname{originalreference}{Original reference}{Original references}
\makeatletter
\newlabel{lemma-ses-statements-coefficients}{{0GJG}{source}{}{sourcepointer-2}{}}
\newlabel{lemma-ses-statements-coefficients@cref}{{[originalreference][0][]0GJG}{[1][0][]source}{}{}{}}
\makeatother
\crefname{originalreference}{original reference}{original references}
\Crefname{originalreference}{Original reference}{Original references}
\makeatletter
\newlabel{lemma-vanishing-easier-coefficients}{{0GJH}{source}{}{sourcepointer-3}{}}
\newlabel{lemma-vanishing-easier-coefficients@cref}{{[originalreference][0][]0GJH}{[1][0][]source}{}{}{}}
\makeatother
\title{Finiteness of the cohomology of a torsion constructible sheaf of modules on a curve over an algebraically closed field}
\author{The Stacks Project Authors}
\date{stacks-project:etale-cohomology.tex:local-snapshot-20261011; GNU FDL 1.2 or later}
\begin{document}
\maketitle
\noindent\textit{Editorial note.} All delivered mathematical text is the author's own, copied from the pinned Stacks Project chapter file etale-cohomology.tex; no statement or proof step was rewritten, reordered or removed.\par\medskip
\noindent\textit{Editorial note.} The counted proof is the author's complete argument: the proof is written out in full: the decomposition into prime-power torsion summands with the additivity input, the reduction to an extension by zero from a dense open, the splitting into irreducible components, and the componentwise application of the companion lemma.\par\medskip
\noindent\textit{Editorial note.} The label theorem-vanishing-affine-curves-coefficients names the counted statement, and the proof environment immediately following it in the source is delivered verbatim as its proof.\par\medskip
\noindent\textit{Editorial note.} The counted statement is carried by the author's own theorem environment in the chapter source.\par\medskip
\noindent\textit{Editorial note.} The statement is an author Theorem environment with two labelled assertions, delivered verbatim; the three cited inputs are kept as external pointers with their canonical Stacks tags.\par\medskip
\noindent\textit{Allowed standard prerequisites (editorial).}\begin{itemize}
\item Additivity of finiteness along short exact sequences. Same chapter, the lemma stating that finiteness of the cohomology of the outer terms of a short exact sequence of torsion sheaves of modules implies finiteness for the middle term. Supplies the reduction from a direct sum decomposition of the coefficient sheaf to its summands.
\item The restriction to a dense open. Same chapter, the lemma comparing the vanishing and finiteness statements for a sheaf and for its extension by zero from a dense open. Supplies the reduction to a sheaf of the form j\_!G over a dense open.
\item The case of an irreducible open. Same chapter, the companion lemma handling the extension by zero along the inclusion of an irreducible open subscheme. Supplies the final step after decomposing the open into irreducible components.
\end{itemize}\medskip
\ifdefined\hypertarget\hypertarget{sourcepointer-1}{}\fi
\noindent\textit{Source reference.} Original reference 0GJE: sources/stacks-etale-cohomology/pinned-original-source.tex, source lines 15927--15927. Prerequisite pointer invoked at proof line(s) 16096; it names a standard fact of the same chapter that the counted argument uses, and every step of the counted proof built on it is retained verbatim in the delivered body.\par\medskip
\ifdefined\hypertarget\hypertarget{sourcepointer-2}{}\fi
\noindent\textit{Source reference.} Original reference 0GJG: sources/stacks-etale-cohomology/pinned-original-source.tex, source lines 15997--15997. Prerequisite pointer invoked at proof line(s) 16083; it names a standard fact of the same chapter that the counted argument uses, and every step of the counted proof built on it is retained verbatim in the delivered body.\par\medskip
\ifdefined\hypertarget\hypertarget{sourcepointer-3}{}\fi
\noindent\textit{Source reference.} Original reference 0GJH: sources/stacks-etale-cohomology/pinned-original-source.tex, source lines 16011--16011. Prerequisite pointer invoked at proof line(s) 16112; it names a standard fact of the same chapter that the counted argument uses, and every step of the counted proof built on it is retained verbatim in the delivered body.\par\medskip
\begin{theorem}
\label{theorem-vanishing-affine-curves-coefficients}
Let $\Lambda$ be a Noetherian ring, let $k$ be an algebraically closed field,
let $X$ be a separated, finite type scheme of dimension $\leq 1$ over $k$,
and let $\mathcal{F}$ be a constructible sheaf of $\Lambda$-modules
on $X_\etale$ which is torsion. Then
\begin{enumerate}
\item
\label{item-finite-prime-to-p-coefficients}
$H^q_\etale(X, \mathcal{F})$ is a finite $\Lambda$-module
if $\mathcal{F}$ is torsion prime to $\text{char}(k)$,
\item
\label{item-finite-proper-coefficients}
$H^q_\etale(X, \mathcal{F})$ is a finite $\Lambda$-module if $X$ is proper.
\end{enumerate}
\end{theorem}

\begin{proof}
without further mention.
Write $\mathcal{F} = \mathcal{F}_1 \oplus \ldots \oplus \mathcal{F}_r$
where $\mathcal{F}_i$ is annihilated by $\ell_i^{n_i}$
for some prime $\ell_i$ and integer $n_i > 0$.
By Lemma \ref{lemma-ses-statements-coefficients} it suffices
to prove the theorem for $\mathcal{F}_i$. Thus we may and do
assume that $\ell^n$ kills $\mathcal{F}$ for some prime $\ell$
and integer $n > 0$.

\medskip\noindent
Since $\mathcal{F}$ is constructible as a sheaf of $\Lambda$-modules,
there is a dense open $U \subset X$ such that $\mathcal{F}|_U$ is a
finite type, locally constant sheaf of $\Lambda$-modules.
Since $\dim(X) \leq 1$ we may assume, after shrinking $U$, that
$U = U_1 \amalg \ldots \amalg U_n$ is a disjoint union of irreducible
schemes (just remove the closed points which lie in the intersections
of $\geq 2$ components of $U$). By
Lemma \ref{lemma-restrict-to-open-coefficients}
we reduce to the case
$\mathcal{F} = j_!\mathcal{G}$ where $\mathcal{G}$ is a finite type,
locally constant sheaf of $\Lambda$-modules on $U$ (and annihilated
by $\ell^n$).

\medskip\noindent
Since we chose $U = U_1 \amalg \ldots \amalg U_n$ with $U_i$ irreducible
we have
$$
j_!\mathcal{G} =
j_{1!}(\mathcal{G}|_{U_1}) \oplus \ldots \oplus
j_{n!}(\mathcal{G}|_{U_n})
$$
where $j_i : U_i \to X$ is the inclusion morphism.
The case of $j_{i!}(\mathcal{G}|_{U_i})$ is handled in
Lemma \ref{lemma-vanishing-easier-coefficients}.
\end{proof}

\end{document}
